% \input{\pSections "sec-redux"}

\section{Redux}

%%    %%    %%    %%    %%    %%    %%    %%    %%    %%
\subsection{Reproduction}

\begin{frame}\frametitle{Comparing Data, Solution Curves}
\center
	%\href{}{
	\begin{overpic}[ scale = 0.9 ]
		{\pLocalGraphics mathematica/Taylor-redux-true}
		\put(38, 51)	{\colorbox{white}{\rd{Original Fit}}}
		\put(60, 40)	{\colorbox{white}{\bl{Least Squares}}}
		\put(60, 35)	{\colorbox{white}{\bl{Minimization}}}
		\put(50, -5)	{\small{\textrm{$t$, s}}}
		\put(-6, 25)	{\small{\textrm{\rotatebox{90}{$R$, cm}}}}
	\end{overpic}
\end{frame}
%
\begin{frame}\frametitle{Seeing the Minimum: Solution Space}
\center
	%\href{}{
	\begin{overpic}[ scale = 0.6 ]
		{\pLocalGraphics mathematica/Taylor-contours-red}
		\put(7, 87)		{\colorbox{white}{$\normt{r(a,b)} = \normt{a t^{b} - R}$}}
		\put(-9, 46)	{\rotatebox{90}{\colorbox{white}{$b$}}}
		\put(49, -5)	{\colorbox{white}{$a$}}
		\put(100, 48)	{\colorbox{white}{\bl{$r_{LS} = 510$}}}
		\put(100, 88)	{\colorbox{white}{\rd{$r =1456$}}}
	\end{overpic}
\end{frame}
%
\begin{frame}\frametitle{Seeing the \href{https://en.wikipedia.org/wiki/Errors_and_residuals}{Residual Error}}
\center
	%\href{}{
	\begin{overpic}[ scale = 0.9 ]
		{\pLocalGraphics mathematica/Taylor-residuals-both-A}
		\put(-6, 25)	{\colorbox{white}{\rotatebox{90}{$error$}}}
		\put(50, -5)	{\colorbox{white}{$t$}}
	\end{overpic}
\end{frame}
\begin{frame}\frametitle{Seminal Papers}
%\caption{default}
\begin{table}[htp]
\begin{center}
	\begin{tabular}{l}
		 %
	\href{run:../local/articles-redux/Taylor50a.pdf}{
	\begin{overpic}[ scale = 0.25 ]
		{\pLocalGraphics taylor/taylor-01.pdf}
		\href{\urlTaylora}{\put(85, 5)	{\colorbox{gray!20}{\cite{taylor1950formationI}}}}
	\end{overpic}}  \\\hline
		 %
	\ \\
	\href{run:../local/articles-redux/Taylor50b.pdf}{
	\begin{overpic}[ scale = 0.25 ]
		{\pLocalGraphics Taylor/taylor-02.pdf}
		\href{\urlTaylorb}{\put(85, 5)	{\colorbox{gray!20}{\cite{taylor1950formationII}}}}
	\end{overpic}} 
		 %
	\end{tabular}
\end{center}
\label{tab:taylor-splash}
\end{table}
\end{frame}
%
\begin{frame}\frametitle{Not \yl{Brook Taylor}}
\begin{table}[htp]
%\caption{default}
\begin{center}
\href{https://en.wikipedia.org/wiki/Brook_Taylor}{
	\begin{tabular}{cc}
		\includegraphics[ width = 1.25in ]{ \pLocalGraphics Taylor/brook-01} &
		\includegraphics[ width = 1.19in ]{ \pLocalGraphics Taylor/brook-02} 
	\end{tabular}
}
\end{center}
%\label{tab:label}
\end{table}
\href{https://mathworld.wolfram.com/TaylorSeries.html}{
\begin{equation*}
	f(x) = f(a) + \frac{x-a}{1!}f' + \frac{(x-a)^{2}}{2!}f''+ \cdots +  \frac{(x-a)^{n}}{n!}f^{(n)} + \cdots
%\label{eq:}
\end{equation*}
}
\end{frame}
%
\begin{frame}\frametitle{Fireball Energy}
\center
	\href{\urlTaylora}{
	\begin{overpic}[ scale = 0.25 ]
		{\pLocalGraphics Taylor/Taylor-01.pdf}
		\put(32, -10)	{\colorbox{white}{$\bl{R = S(\gamma) \sqrt[5]{\frac{E\, t^{2}}{\rho_{0}} }}$}}
		\put(65, -20)	{\colorbox{gray!20}{\cite[p. 172]{taylor1950formationI}}}
	\end{overpic}}
\end{frame}
%
\begin{frame}\frametitle{Fireball Photos: Trinity}
\center
		%
	\href{https://homepage.physics.uiowa.edu/~ghowes/teach/phys7729/docs/Taylor50b.pdf\#page=9}{
	\begin{overpic}[ scale = 0.175 ]
		{\pLocalGraphics taylor/p-photos-01}
		%\put(50, 50)	{\colorbox{white}{$a+b$}}
	\end{overpic}} 
		%
	\href{https://homepage.physics.uiowa.edu/~ghowes/teach/phys7729/docs/Taylor50b.pdf\#page=2}{
	\begin{overpic}[ scale = 0.385 ]
		{\pLocalGraphics taylor/p-table-1}
		%\put(50, 50)	{\colorbox{white}{$a+b$}}
	\end{overpic}}
		%
\end{frame}

\begin{frame}\frametitle{\href{https://homepage.physics.uiowa.edu/~ghowes/teach/phys7729/docs/Taylor50b.pdf\#page=3}{Taylor's Original Figure 1}}
\center
	\href{https://homepage.physics.uiowa.edu/~ghowes/teach/phys7729/docs/Taylor50b.pdf\#page=3}{
	\begin{overpic}[ scale = 0.425 ]
		{\pLocalGraphics taylor/fig-01}
		\pause
		\put(40, 40)	{\colorbox{white}{Method of Choice: \bl{Least Squares}}}
	\end{overpic}}
\end{frame}

%
\begin{frame}\frametitle{Validity of Similarity Relationship}
\center
	\href{https://homepage.physics.uiowa.edu/~ghowes/teach/phys7729/docs/Taylor50b.pdf\#page=3}{
	\begin{overpic}[ scale = 0.5 ]
		{\pLocalGraphics taylor/Taylor-basic}
		\put(-45, 105)	{\colorbox{white}{$\frac{5}{2}\log R(t) = (\bl{\solnb} \pm \bl{\errb})\log t + (\solna \pm \erra)$}}
		\put(-15, 40)	{\rotatebox{90}{\colorbox{white}{$\frac{5}{2}\log R$}}}
		\put(37, -8)	{\colorbox{white}{$\log t$}}
	\end{overpic}}
\end{frame}
%
\begin{frame}\frametitle{Least Squares Problem}
	\huge{
	\begin{equation}
		(\bl{a}, \bl{b})_{LS}\equiv \argmin\limits_{(\bl{a},\bl{b})\in\mathbb{R}^{2}} \normt{\bl{a} t^{\bl{b}} - R}
	\label{eq:solution definition}
	\end{equation}
	}
\end{frame}

\begin{frame}\frametitle{Solution Using \href{https://www.wolfram.com/mathematica/	}{Mathematica}}
\center
	\href{run:../local/articles-redux/2023-film-working-group.pdf}{
	\begin{overpic}[ scale = 0.375 ]
		{\pLocalGraphics mathematica/norm-min}
		\pause
		\put(55, 5)	{\colorbox{white}{$(\bl{a}, \bl{b})_{LS}\equiv \argmin\limits_{(\bl{a},\bl{b})\in\mathbb{R}^{2}} \normt{\bl{a} t^{\bl{b}} - R}$}}
	\end{overpic}}
\end{frame}


\begin{frame}\frametitle{Numeric Alert}
\center
	%\href{}{
	\begin{overpic}[ scale = 0.375 ]
		{\pLocalGraphics mathematica/ouch}
		%\put(50, 50)	{\colorbox{white}{$a+b$}}
	\end{overpic}
\end{frame}

%
\begin{frame}\frametitle{Error Analysis}
\begin{equation}
	\sigma_{R} = b \sigma_{t} \frac{R}{t}
\label{eq:error-powers}
\end{equation}
\end{frame}
%

%\begin{frame}\frametitle{Beirut Yield Estimates: Academia}
%\begin{enumerate}
%	\item \href{https://link-springer-com.libproxy.unm.edu/article/10.1007/s00193-021-01031-9}
%		{\fullcite{aouad2021beirut}}
%	\item \href{https://arxiv.org/pdf/2009.05674.pdf}
%		{\fullcite{diaz2021explosion}}
%	\item \href{https://aiche-onlinelibrary-wiley-com.libproxy.unm.edu/doi/epdf/10.1002/prs.12203}
%		{\fullcite{pasman2020beirut}}
%\end{enumerate}
%\end{frame}
%%    %%    %%    %%    %%    %%    %%    %%    %%    %%
\subsection{Rudiments}

%	\input{\pSections "slide-ftolathm"}

\begin{frame}\frametitle{\lsq Untangles The Null Spaces}
	%
\begin{figure}[t]
\centering
	\begin{overpic}[ scale = 0.75 ]{\pLocalGraphics bev/colored-masks-02}
			%
	\end{overpic}
\label{fig:Seeing the subspace decomposition}
\caption{Vectors in \href{https://mathworld.wolfram.com/ColumnSpace.html}{\bl{range spaces} $\left( \brnga{} \right.$ or $\left. \brnga{*}\right)$} are colored in blue. Vectors in \href{https://mathworld.wolfram.com/NullSpace.html}{\rd{null spaces} $\paren{\rnlla{*} \text{ or } \rnlla{}}$} are colored in red. Amalgam vectors with components in both spaces are the traditional black.}
\end{figure}
	%
\end{frame}

	\input{\pSections "ssec-repro-classic"}
	\input{\pSections "ssec-repro-non-linear"}	
	\input{\pSections "ssec-repro-mcaninch"}


\endinput  %  ==  ==  ==  ==  ==  ==  ==  ==  ==
